% Supplementary material: moved from Sections 7.6 and 7.7 of the main text.

\section{Conversion Pipeline and Barriers to Adoption}
\label{supp:agentic}

\subsection{Converting an LLM Agent to an SLM Agent}
\label{subsec:agentic_conversion}

For teams with an existing LLM-only agent, \citet{belcak2025slm_agentic} propose a six-step conversion pipeline, which we summarize here because it operationalizes the position; the steps are theirs. \textbf{S1, instrumentation:} log every non-interface model call (inputs, outputs, latency) from the running agent, with encryption and access control. \textbf{S2, curation:} strip personally identifiable and protected health information and application secrets from the logs. \textbf{S3, clustering:} apply unsupervised clustering to the curated calls to surface the recurring request types, which are the specialization targets. \textbf{S4, selection:} for each cluster, choose a candidate \gls{slm} by matching required capability against benchmarks, license, and deployment constraints. \textbf{S5, specialization:} fine-tune the chosen \gls{slm} on the cluster's curated data, typically with \gls{lora} or \gls{qlora}. \textbf{S6, iteration:} re-run the loop periodically as the traffic distribution drifts. Two additions follow from H. Step~S3 is where the conditions are assessed: each cluster's output schema (H1), input length (H3), and the availability of validity and correctness checks (H4a, H4b) can be checked by inspection, and its volume can be compared with a pilot estimate of the break-even $N^{*}$ (H2), recorded before the comparison; clusters that fail are left with the generalist. And each converted cluster should be evaluated on cost per completed task against the unconverted agent, which is the measurement the literature lacks.

\subsection{Barriers to Adoption}
\label{subsec:agentic_barriers}

\citet{belcak2025slm_agentic} identify three reasons the shift has lagged the technical case. First, large centralized-inference investments create sunk-cost inertia and an operational bias toward one big model behind one API (B1). Second, \glspl{slm} are still evaluated on generalist leaderboards (\gls{mmlu}, broad chat benchmarks) that are misaligned with agentic utility, so their fitness for delegated sub-tasks is under-measured (B2). Section~\ref{sec:reported} shows that the generalist record is also too inconsistent across vendors to support cross-model comparison, so it cannot be used to choose among small models even for general capability. Third, \gls{slm} releases receive far less attention than frontier launches, so awareness lags capability (B3). The first and third are not technical obstacles. The second is different, and Section~\ref{subsec:agentic_evidence} shows it to be partly a real evidential gap rather than only a misdirected emphasis: agent-specific evaluation exists and is maturing \citep{patil2025bfcl,yao2024taubench,barres2025tau2,liu2024agentbench,kapoor2025hal}, and developer reports support per-call competence for specialized small models, but independent multi-step evaluations of current sub-10B models are sparse, and no study compares cost per completed task for small-model-first and large-model agents. Until that changes, the strength of the case for hypothesis~H will remain uneven in the way Table~\ref{tab:agentic_evidence} records.

\subsection{Learned Routing and Cascades}
\label{subsubsec:agentic_routing}

Three designs recur. \emph{Cascades} try the cheap model first and escalate on a quality signal: FrugalGPT \citep{chen2023frugalgpt} learns which sequence of \gls{llm} APIs to query, and AutoMix \citep{aggarwal2024automix} escalates when the small model's own few-shot self-verification judges its answer unreliable, which is the pattern H prescribes when the checker of H4 is the model itself. \emph{Predictive routers} decide before generation: Hybrid~LLM \citep{ding2024hybridllm} routes by predicted query difficulty with a quality level tunable at test time, and RouteLLM \citep{ong2024routellm} trains a lightweight router on human-preference data to predict whether a weak model's answer would be judged as good as a strong model's, reporting transfer of the trained router to unseen strong/weak model pairs. \emph{Decomposing routers} push the decision below the query level: Division-of-Thoughts \citep{division_of_thoughts_2025} decomposes a task into sub-questions and assigns each to a local \gls{slm} or a cloud \gls{llm} according to estimated difficulty, for on-device agents. The open problems are difficulty estimation without an oracle, routing objectives that jointly weigh accuracy, latency, privacy sensitivity, and network availability, and handoff that does not produce a visible seam in output quality \citep{demystifying2025slm}.

A different form of SLM--LLM collaboration leaves the large model in charge of every output and uses the small one only to make it faster. In speculative decoding a small draft model proposes several tokens and the large model verifies them in one forward pass, accepting a prefix, so that the output distribution is that of the large model \citep{leviathan2023speculative,chen2023specsampling}; Medusa \citep{cai2024medusa} and EAGLE \citep{li2024eagle} replace the separate draft model with lightweight heads on the large model itself. This is not delegation in the sense of H, since no call's answer comes from the small model, but it is the most widely deployed way in which small models already reduce the cost of large-model inference, and it is complementary to routing: an escalated call can itself be served speculatively.

